\documentclass[12pt]{article}
\usepackage{../../includes/lecture_notes}
\usepackage{../../includes/math}
\usepackage{../../includes/uark_colors}

\title{Regression}
\author{Prof. Kyle Butts}
\date{}

\hypersetup{
  colorlinks = true,
  allcolors = ozark_mountains,
  breaklinks = true,
  bookmarksopen = true
}

\begin{document}
\maketitle
\begin{multicols}{2}

  \section*{Conditional Expectation Function}

  The \textbf{conditional expectation function (CEF)} is one of the most fundamental concepts in applied econometrics.
  At its core, it asks a simple question: given that we know some characteristics about a unit, what do we expect their outcome to be?

  Formally, we define the CEF as:
  $$m(\bm{x}) \equiv \mathbb{E}[y_i|\bm{X}_i = \bm{x}]$$

  This reads as ``$m(\bm{x})$ is the expected value of $y_i$ conditional on the unit having characteristics $\bm{X}_i = \bm{x}$.''
  The beauty of this definition is that it connects directly to how we naturally think about prediction and relationships in data.

  The most straightforward way to estimate the CEF for any given value $\bm{x}$ is remarkably simple: find all units in your data with $\bm{X}_i = \bm{x}$ and take their average outcome.
  This is the simplest estimation procedure for the CEF.
  For each possible value of $\bm{x}$,
  \begin{enumerate}
    \item Subset your data to units with $\bm{X}_i = \bm{x}$
    \item Take average of $y_i$ for those units. Call $m(\bm{x})$ this subsample average.
  \end{enumerate}

  Let's start with the simplest case: a single discrete variable.
  Consider wages ($w_i$) and college attendance ($D_i$), where $D_i = 1$ indicates college attendance and $D_i = 0$ indicates no college.
  The CEF becomes:
  $$
  m(1) = \mathbb{E}[w_i|D_i = 1] \quad\text{and}\quad m(0) = \mathbb{E}[w_i|D_i = 0]
  $$

  This is exactly what we'd compute: the average wage for college attendees and non-attendees.

  Now suppose we have two discrete variables: college attendance ($D_i$) and gender, where $F_i = 1$ indicates female.
  We now have four distinct combinations: $(D_i, F_i) \in \left\{ (0,0), (1,0), (0,1), (1,1) \right\}$.
  Estimating the CEF requires computing four separate averages for each group.
  This highlights an important point: to fully characterize the relationship, we need to account for all possible interactions between our variables.

  Moving to continuous variables introduces new challenges.
  If $X_i$ is continuous, we can no longer find units with exactly $X_i = x$.
  Instead, we must find units with $X_i \approx x$.
  But how close is ``close enough''?
  The standard approach is to create many small \textbf{bins} across the range of $X$ and average outcomes within each bin.
  This non-parametric estimation approach works well when we have plenty of data and few variables.

  However, as we add more variables, we quickly encounter the \textbf{curse of dimensionality}.
  With multiple continuous variables, the number of distinct combinations of $\bm{x}$ values grows exponentially.
  In practice, most combinations will have few or no observations, making reliable estimation impossible.
  This fundamental limitation pushes us toward modeling the CEF rather than estimating it directly.

  \section*{Modelling the CEF}

  The curse of dimensionality forces us to make assumptions about the functional form of $m(\bm{x})$.
  Rather than estimating the CEF separately for every possible combination of characteristics, we must specify a model that can leverage all our data efficiently.

  The most common approach is the \textbf{linear model}:
  $$
  m(\bm{x}) = \bm{x}' \beta = \sum_{k=1}^p x_{k} \beta_k
  $$
  We assume that one of the $x_k = 1$ for all units, so that we have an intercept.
  This assumes that the expected outcome changes linearly with each variable $x_k$.

  But linearity often fails to capture important features of real relationships.
  Consider a few examples where constant marginal effects seem implausible:
  \begin{itemize}
    \item Returns to experience in the labor market typically exhibit diminishing returns.
      The 20th year of experience adds less to wages than the second year.

    \item In housing markets, an additional square foot adds less value to a mansion than to a studio apartment.

    \item The return to experience may depend critically on whether someone already has a college degree.
      Here, there is complementarity between different types of human capital.
  \end{itemize}

  These examples suggest we need more flexibility while maintaining the tractability of linear models.
  The key insight is to distinguish between the variables we observe, $\bm{X}_i$, and the \textbf{terms} we include in our model.
  In general, we write
  $$
  \bm{W}_i = \bm{g}(\bm{X}_i) = \left( g_1(\bm{X}_i), \dots, g_K(\bm{X}_i) \right)'.
  $$
  For diminishing returns, we might include polynomial terms: if $X_i$ is years of experience, we could include both $X_i$ and $X_i^2$ in $\bm{W}_i$.
  For complementarity, we include interaction terms: the return to experience might be $\beta_1 X_{1i} + \beta_2 X_{1i} \cdot D_i$ where $D_i$ indicates college education.
  For highly flexible relationships, we can use bins: divide a continuous variable into ranges and include indicator variables for each range.

  The selection of terms to include in a linear regression model is a bit of an art.
  Including too many terms makes the model at risk of overfitting and makes the output more difficult to interpret.
  Too few terms risks making strong restrictions on how $\bm{X}_i$ and $y_i$ relate.

  One very common setting is wanting to consider a highly flexible relationship of one key variable ($X_{1,i}$) while maintaining a simple linear control for other variables ($\bm{Z}_i$).
  This model,
  $$
  y_i = \mu(X_{1,i}) + \bm{Z}_i' \beta + u_i,
  $$
  is often called the \textbf{partially linear model} because the function we let $\mu(X_{1,i})$ be very flexible while having a simple linear model for the other covariates.

  There are many ways to estimate this model, but one popular method is called Binscatter Regression and is implementable via the \texttt{binsreg} package.

  \section*{Ordinary Least Squares}

  After selecting a model for the conditional expectation function, i.e. a choice of transformations of $\bm{X}_i$ in $\bm{W}_i$, we move to estimating this model.
  Our model is given by:
  $$
  y_i = \bm{W}_i' \beta + u_i.
  $$
  At a high level, our goal is to use a sample of observations to learn which values of $\beta$ best predict the outcome.
  But what constitutes ``best''?
  There are many choices on how to evaluate model fit, which we call the \textbf{loss function}.

  By far the most common choice is the mean-squared prediction error:
  $$
  \text{MSPE}(\beta) \equiv \sum_{i=1}^n (y_i - \bm{W}_i' \beta)^2.
  $$

  The \textbf{Ordinary Least Squares} (OLS) estimator is the value of $\beta$ that minimizes MSPE:
  $$
  \hat{\beta}_{\text{OLS}} \equiv \argmin_\beta \sum_{i=1}^n (y_i - \bm{W}_i' \beta)^2.
  $$

  \section*{Matrix Notation}

  Let $\bm{Y}$ be the $n \times 1$ vector of $y_i$. Let $\bm{W}$ be the $n \times K$ matrix stacking $\bm{W}_i'$:
  $$
  \bm{W} =
  \begin{bmatrix}
    \bm{W}_1' \\
    \vdots    \\
    \bm{W}_n'
  \end{bmatrix}.
  $$
  We generally \emph{always} assume you have an intercept, i.e. $W_{i1} = 1$.
  Our model in matrix form is:
  $$
  \bm{Y} = \bm{W} \beta + \bm{u}.
  $$

  \section*{Deriving the OLS Estimator}

  We want to minimize the sum of squared residuals:
  $$
  \hat{\beta}_{\text{OLS}} = \argmin_\beta (\bm{Y} - \bm{W} \beta)'(\bm{Y} - \bm{W} \beta).
  $$

  Expanding this product:
  \begin{align*}
    &(\bm{Y} - \bm{W} \beta)'(\bm{Y} - \bm{W} \beta) \\
    &\quad =
    \bm{Y}' \bm{Y} - \beta' \bm{W}' \bm{Y} - \bm{Y}' \bm{W} \beta + \beta' \bm{W}' \bm{W} \beta.
  \end{align*}

  Squint your eyes and this looks like a quadratic function.
  This is indeed true, in the linear algebra world, this is what a quadratic function looks like.
  Taking the derivative and setting it equal to zero yields the first-order condition:
  \begin{align*}
    \frac{\partial}{\partial \beta} \left( \bm{Y}' \bm{Y} - \beta' \bm{W}' \bm{Y} - \bm{Y}' \bm{W} \beta + \beta' \bm{W}' \bm{W} \beta \right) = 0.
  \end{align*}

  Using the rules of matrix derivatives:
  \begin{align*}
    0 - \bm{W}' \bm{Y} - \bm{W}' \bm{Y} + 2 \bm{W}' \bm{W} \beta = 0.
  \end{align*}

  Setting this equal to zero gives:
  $$
  \left(\bm{W}' \bm{W} \right) \hat{\beta}_{\text{OLS}} = \bm{W}' \bm{Y}.
  $$

  Solving for $\beta$:
  $$
  \hat{\beta}_{\text{OLS}} = \left(\bm{W}' \bm{W} \right)^{-1} \bm{W}' \bm{Y}.
  $$

  \section*{Intuition for the OLS Formula}

  Let's check that this formula makes sense in simple cases.

  \noindent\textbf{Intercept-only model:} Say we have just an intercept ($W_{i1} = 1$), so $\bm{W} = \iota$. Then:
  \begin{itemize}
    \item $\bm{W}' \bm{W} = \iota' \iota = n$.
    \item $\bm{W}' \bm{Y} = \iota' \bm{Y} = \sum_{i=1}^n Y_i$.
  \end{itemize}
  Consequently, $\hat{\beta}_{\text{OLS}} = \frac{1}{n} \sum_{i=1}^n Y_i$, the sample mean.

  \noindent\textbf{Bivariate regression:} Consider an intercept and a single explanatory variable $W_{i2}$.
  By the Frisch-Waugh-Lovell theorem, this is equivalent to regressing $Y_i - \bar{Y}$ on $W_{i2} - \bar{W}_2$.
  In this case:
  \begin{itemize}
    \item $\bm{W}' \bm{W} = \sum_i (W_{i2} - \bar{W}_2)^2$ is $(n-1)$ times the sample variance of $W_{i2}$.
    \item $\bm{W}' \bm{Y} = \sum_i (W_{i2} - \bar{W}_2)(Y_i - \bar{Y})$ is $(n-1)$ times the sample covariance.
  \end{itemize}

  Consequently, we have the bivariate regression formula: $\hat{\beta}_{\text{OLS}} = \widehat{\cov}(W_{i2}, Y_i) / \widehat{\var}(W_{i2})$.

  \noindent\textbf{Multivariate regression:} More generally, when we have $K-1$ covariates and an intercept, the regression is equivalent to one where $Y$ and all covariates are demeaned (without an intercept).
  Then:
  $$
  \hat{\beta}_{\text{OLS}} = \left[ \widehat{\var}(\bm{W}_i) \right]^{-1} \widehat{\cov}(\bm{W}_i, Y_i).
  $$

  \section*{Statistical Properties of OLS}

  In repeated sampling, we will get different sample of observations.
  This will create different estimates of $\hat{\beta}$.

  Say the we have \textbf{correctly specified} the conditional expectation function, i.e. the true model $y_i = \bm{W}_i' \beta_0 + u_i$.
  Then we can assume $\expec{u_i}{\bm{W}_i} = 0$.

  \bigskip
  \emph{\textbf{Side Note:}} All of the below derivations are still true without assuming we correctly specify the model.
  In that case, $\beta_0$ corresponds to the coefficient we would get if we ran the regression on the full population (the population least squares coefficient).
  This would imply $\text{Cov}(u_i, \bm{W}_i) = 0$, i.e. there's no linear covariance between the error and $\bm{W}$.

  Plugging this into our OLS estimator:
  \begin{align*}
    \hat{\beta}_{\text{OLS}}
    & = \left(\bm{W}' \bm{W} \right)^{-1} \bm{W}' \bm{Y}                                                            \\
    & = \left(\bm{W}' \bm{W} \right)^{-1} \bm{W}' \left( \bm{W} \beta_0 + \bm{u} \right)                            \\
    & = \left(\bm{W}' \bm{W} \right)^{-1} \bm{W}' \bm{W} \beta_0 + \left(\bm{W}' \bm{W} \right)^{-1} \bm{W}' \bm{u} \\
    & = \beta_0 + \left(\bm{W}' \bm{W} \right)^{-1} \bm{W}' \bm{u}.
  \end{align*}

  This decomposition shows that our estimator equals the true population parameter plus a noise term that depends on the random errors.

  \section*{Unbiasedness}

  Taking expectations of the expression above and using $\expec{u_i}{\bm{W}_i} = 0$ (or the weaker 0 covariance):
  \begin{align*}
    \expec{\hat{\beta}_{\text{OLS}}}
    & = \beta_0 + \expec{ \left(\bm{W}' \bm{W} \right)^{-1} \bm{W}' \bm{u} } \\
    & = \beta_0.
  \end{align*}

  On average across repeated samples, our estimator equals the true population parameter.
  This means our coefficient estimates are \textbf{unbiased} for the population values.

  \section*{The Sample Distribution of OLS Coefficients}

  Subtracting $\beta_0$ and multiplying by $\sqrt{n}$:
  $$
  \sqrt{n} \left( \hat{\beta}_{\text{OLS}} - \beta_0 \right) =
  \left( \frac{1}{n} \bm{W}' \bm{W} \right)^{-1} \frac{1}{\sqrt{n}} \sum_{i=1}^n \bm{W}_i u_i.
  $$

  The term $\frac{1}{\sqrt{n}} \sum_{i=1}^n \bm{W}_i u_i$ has mean zero (from unbiasedness) and variance $\expec{\bm{W}' \bm{u} \bm{u}' \bm{W}} = \bm{W}' \Sigma \bm{W}$ where $\Sigma = \expec{\bm{u} \bm{u}'}$.

  Applying the Central Limit Theorem, we have that $\frac{1}{\sqrt{n}} \bm{W}' \bm{u}$ is approximately normally distributed. Therefore:
  $$
  \hat{\beta}_{\text{OLS}} \sim
  \mathcal{N}\left( \beta_0, \left(\bm{W}' \bm{W} \right)^{-1} \bm{W}' \Sigma \bm{W} \left(\bm{W}' \bm{W} \right)^{-1} \right).
  $$

  This is a remarkable result. Even if the underlying data are not normally distributed, the OLS estimator has an approximately normal distribution in large samples.

  \section*{Variance and Standard Errors}

  The variance is a $K \times K$ matrix. The diagonal elements $\var{\hat{\beta}_{\text{OLS}, k}}$ give the variance of each coefficient estimate. Take the square root to get the standard deviation of $\hat{\beta}_{\text{OLS}, k}$.

  The off-diagonal elements tell us how slope coefficients are correlated with one another across repeated samples. This matters when we want to test multiple hypotheses simultaneously.

  \bigskip
  The main challenge of estimating the variance-covaraince matrix entails the matrix $\Sigma$ which is a $N \times N$ matrix.
  The element $\Sigma_{i,j}$ denotes the covariance between the error terms for observation $i$ and $j$.

  \section*{Heteroskedasticity-Robust Standard Errors}
  By default, software will produce `IID standard errors'.
  Do not use these. Instead, we will use (at minimum) `HC1' standard errors.
  These are often called `robust standard errors', but I prefer the full name heteroskedasticity-robust standard errors.

  This imposes that $\Sigma_{i,j} = 0$ whenever $i \neq j$, but lets the variance of the error term vary by individual (e.g. wage shocks to high-skilled workers might be larger than low-skilled workers).

  We estimate the variance using residuals $\hat{u}_i = y_i - \bm{W}_i' \hat{\beta}_{\text{OLS}}$. The HC1 variance estimator is:
  $$
  \widehat{\var}(\hat{\beta}_{\text{OLS}}) = \left(\bm{W}' \bm{W} \right)^{-1} \bm{W}' \hat{\Sigma} \bm{W} \left(\bm{W}' \bm{W} \right)^{-1},
  $$
  where $\hat{\Sigma}$ has $\hat{u}_i^2$ on the diagonal.

  Taking the square root of the $k$-th diagonal element gives the \textbf{standard error} of $\hat{\beta}_{\text{OLS}, k}$: our estimate of the standard deviation of the sampling distribution.

  \section*{Cluster-Robust Standard Errors}

  Previously we discussed ``robust'' or ``HC1'' standard errors, which assume \emph{independent} errors: a shock to one unit is unrelated to all other units.

  In many settings, independent errors are implausible because unit-level shocks are correlated. Think about:
  \begin{itemize}
    \item Students in the same school.
    \item Workers in the same labor market.
    \item Patients at the same hospital.
  \end{itemize}

  When errors are correlated across units, our ``robust'' standard errors are typically too small. The intuition: correlated errors mean we have fewer ``effective observations'' than $n$.

  \section*{Clustering}

  The most common solution is to model errors as being \emph{clustered}:
  \begin{itemize}
    \item Units within the same cluster are allowed to have arbitrarily correlated errors.
    \item Units in different clusters have zero correlation.
  \end{itemize}

  Let $g_i$ denote the cluster that unit $i$ belongs to.
  Then $\sigma_{i,j} = 0$ when $g_i \neq g_j$, but $\sigma_{i,j}$ is unrestricted when $g_i = g_j$.

  The cluster-robust variance estimator uses the block-diagonal structure of the covariance matrix. We estimate each $\sigma_{i,j}$ as the product of regression residuals $\hat{u}_i \hat{u}_j$.

  The intuition for why standard errors are larger: we have fewer ``distinct'' observations when errors can be correlated within groups.
  In the worst case (extreme correlation), you only have $G$ observations (one per cluster).
  In the best case (independence within cluster), you have $n$.

  The rule of thumb is you want at least 20-30 clusters, each with multiple observations.
  If you have fewer groups, use the wild cluster bootstrap.

  \section*{Fixed Effects vs. Clustering}

  We could include cluster fixed effects:
  $$y_i = \alpha_{g_i} + \beta x_i + \epsilon_i,$$
  where $\alpha_{g_i}$ is the common effect for all units in the same group.

  But why would we still need clustered standard errors?

  Fixed effects assume \emph{all units} in the group are affected \emph{by the same amount}. But there can be more general forms of correlated errors within a group. Some people might be more impacted by shocks than others. Clustered standard errors can still be necessary even when using cluster fixed effects.

  \section*{Confidence Intervals}

  Since $\hat{\beta}_{\text{OLS}, k}$ has an approximately normal sampling distribution, we can form confidence intervals just as before:
  $$
  \left[ \hat{\beta}_{k} - 1.96 \cdot \text{SE}(\hat{\beta}_{k}), \ \hat{\beta}_{k} + 1.96 \cdot \text{SE}(\hat{\beta}_{k}) \right]
  $$

  Once you choose an estimator for the variance-covariance matrix of the OLS coefficients, then we take the square-root of the diagonal elements to form $\text{SE}(\hat{\beta}_{\text{OLS}, k})$.
  The interpretation is the same as before: across repeated samples, 95\% of those confidence intervals will contain the true value $\beta_{0,k}$.

  \section*{Forecasting with the Regression Model}

  We have a fitted model $\hat{y} = \bm{W}' \hat{\beta}_{\text{OLS}}$ with:
  $$
  \hat{\beta}_{\text{OLS}} \sim \mathcal{N}\left( \beta_0, \bm{V} \right),
  $$
  where $\bm{V} = \left(\bm{W}' \bm{W} \right)^{-1} \bm{W}' \Sigma \bm{W} \left(\bm{W}' \bm{W} \right)^{-1}$.

  Now suppose we want to evaluate this model at a particular value of the covariates, call it $\bm{w}$. The forecasted value is:
  $$
  \hat{Y} = \bm{w}' \hat{\beta}_{\text{OLS}} = \sum_{k=1}^K w_k \hat{\beta}_{\text{OLS}, k}.
  $$

  Uncertainty from our regression coefficients translates to uncertainty about $\hat{Y}$. We can write:
  \begin{align*}
    \hat{Y} & = \bm{w}' \hat{\beta}_{\text{OLS}}                                                                                                       \\
    & = \underbrace{\bm{w}' \beta_0}_{= f_0(\bm{w}) = \expec{Y}{\bm{X} = \bm{w}}} + \bm{w}' \left( \hat{\beta}_{\text{OLS}} - \beta_0 \right).
  \end{align*}

  The forecasted value is the conditional expectation function (assuming our model is correct) plus sampling noise.

  \section*{Inference on Forecasts}

  Since $\hat{\beta}_{\text{OLS}}$ is normally distributed and we multiply it by a fixed vector $\bm{w}$, the forecast is also normally distributed:
  $$
  \hat{Y} = \bm{w}' \hat{\beta}_{\text{OLS}} \sim \mathcal{N}\left( \bm{w}' \beta_0, \ \bm{w}' \bm{V} \bm{w} \right).
  $$

  This is exactly what we covered earlier: forecasts are functions of the data, so they have a sample distribution.
  The Central Limit Theorem ensures this distribution is approximately normal.
  And from the normal distribution, we can form confidence intervals.

  The standard error of the forecast is $\sqrt{\bm{w}' \bm{V} \bm{w}}$, and the 95\% confidence interval is:
  $$
  \hat{Y} \pm 1.96 \cdot \sqrt{\bm{w}' \bm{V} \bm{w}}.
  $$

  \section*{Marginal (Predictive) Effects}

  Often we want to compare forecasted values at two different points: $\bm{w}_1$ and $\bm{w}_2$.

  The simplest approach is to compare $\bm{w}_1' \hat{\beta}_{\text{OLS}}$ and $\bm{w}_2' \hat{\beta}_{\text{OLS}}$ directly.

  It is important to remember that the goal of forecasting is to predict $Y$ as well as possible. When units have larger $x_1$, maybe they tend to have larger $x_2$ as well. Regression will use that information when forming predictions.

  Avoid making causal claims that \emph{changing} $w_k$ \emph{causes} a change in $Y$. Instead, use language like:

  \textbf{Correct:} ``Our regression model predicts that a one unit increase in $w_k$ is associated with a $\hat{\beta}_{\text{OLS},k}$ unit change in $Y$.''

  \textbf{Incorrect:} ``Increasing $w_k$ by one unit increases $Y$ by $\hat{\beta}_{\text{OLS},k}$ units.''

  \section*{Derivatives as Marginal Effects}

  Sometimes we want to think about changing $X_\ell$ instead of changing $W_\ell$ directly. For example, if we change age, we change both age and age$^2$ in our model.

  To make this precise, write our model as:
  $$
  \hat{f}(\bm{X}) = \sum_{k=1}^K g_k(\bm{X}) \hat{\beta}_{\text{OLS}, k}.
  $$

  We can ask how $\hat{f}(\bm{X})$ changes when we change one element of $\bm{X}$, say $x_\ell$. Taking the derivative:
  $$
  \frac{\partial}{\partial x_\ell} \hat{f}(\bm{X}) = \sum_{k=1}^K \frac{\partial}{\partial x_\ell} g_k(\bm{X}) \hat{\beta}_{\text{OLS}, k}.
  $$

  This is called the \textbf{marginal (predictive) effect} of $x_\ell$. Note that I say ``predictive'' rather than ``causal'' to emphasize this is not the effect of experimentally changing $x_\ell$ for a unit.

  This formula holds fixed all other variables at their original covariate values and only changes $x_\ell$. Note that multiple $g_k$ terms can change when we change a single $x_\ell$.

  In the special case where we include each variable linearly ($g_k(\bm{X}) = X_k$), this reduces to the standard $\hat{\beta}_{\text{OLS}, k}$ being our estimated marginal effect.

  For example, if we have a quadratic model $\hat{w}_i = \hat{\beta}_0 + \hat{\beta}_1 \text{age}_i + \hat{\beta}_2 \text{age}_i^2$, then:
  $$
  \frac{\partial}{\partial \text{age}} \hat{w}_i = \hat{\beta}_1 + 2\hat{\beta}_2 \text{age}_i.
  $$

  The marginal effect of age depends on the person's age itself. At age 25, the effect is $\hat{\beta}_1 + 50\hat{\beta}_2$. At age 55, it is $\hat{\beta}_1 + 110\hat{\beta}_2$.
\end{multicols}
\end{document}
